\chapter{Fundamentação Teórica}
\label{cap:fundamentacao}

Este capítulo reúne os fundamentos que sustentam o trabalho, organizados do geral para o
específico, conforme a natureza da contribuição proposta. Parte-se das boas práticas de
Engenharia de Software aplicáveis a sistemas embarcados e robóticos
(Seção~\ref{sec:es-embarcados}), passando pela refatoração
(Seção~\ref{sec:refatoracao}), pela documentação técnica (Seção~\ref{sec:documentacao}),
pelo versionamento e organização de repositórios (Seção~\ref{sec:versionamento}) e pela
arquitetura baseada em componentes materializada no ROS (Seção~\ref{sec:ros}). Por fim,
delimitam-se os fundamentos de domínio (Seção~\ref{sec:dominio}) --- reabilitação,
realidade virtual, eletroestimulação e sensoriamento inercial --- na medida necessária
para compreender os dois estudos de caso.

\section{Engenharia de Software para Sistemas Embarcados e Robóticos}
\label{sec:es-embarcados}

A Engenharia de Software é a disciplina que se ocupa da produção de software de qualidade,
dentro de prazos e custos previsíveis, ao longo de todo o seu ciclo de vida
\cite{sommerville2011}. Um princípio central da disciplina é o de que o custo dominante de
um sistema não está na sua construção inicial, mas na sua \emph{evolução} --- as
sucessivas alterações necessárias para corrigir defeitos, adaptar-se a novos ambientes e
incorporar requisitos. Atributos como manutenibilidade, modularidade e baixo acoplamento,
portanto, não são refinamentos estéticos, mas determinantes diretos do custo e da
viabilidade de manter um sistema vivo ao longo do tempo.

Esses princípios assumem contornos particulares quando o software é \emph{embarcado} e
\emph{robótico}, como é o caso dos sistemas do Projeto EMA. Nesse domínio, o software não
existe isoladamente: ele controla e é controlado por hardware --- sensores, atuadores,
microcontroladores --- sujeito a restrições de tempo real, de segurança e de recursos
computacionais limitados. Um comando enviado a um estimulador elétrico ou a leitura de um
sensor inercial têm consequências físicas imediatas, o que eleva a criticidade de falhas.
Ao mesmo tempo, sistemas robóticos são intrinsecamente distribuídos e heterogêneos, o que
torna a modularidade e a padronização das interfaces entre componentes uma necessidade, e
não uma opção.

Um agravante recorrente em laboratórios de pesquisa é a tensão entre a cultura de
prototipagem rápida --- orientada a demonstrar viabilidade e produzir resultados
publicáveis --- e as boas práticas de engenharia que garantem a longevidade do software. O
protótipo que ``funciona uma vez'' cumpre o seu papel científico imediato, mas, na ausência
de documentação, testes e organização, converte-se em um passivo técnico
(\emph{technical debt}) que onera todas as gerações seguintes. As seções a seguir
apresentam o instrumental que este trabalho mobiliza para diagnosticar e reduzir esse
passivo: refatoração, documentação, versionamento e arquitetura baseada em componentes.

\section{Refatoração}
\label{sec:refatoracao}

Refatoração é o processo de alterar a estrutura interna de um software sem modificar o seu
comportamento externo observável, com o objetivo de torná-lo mais simples, legível e
manutenível \cite{fowler2018}. Trata-se de uma disciplina de melhoria contínua: em vez de
uma reescrita monolítica e arriscada, aplicam-se transformações pequenas e verificáveis que
preservam a funcionalidade enquanto aprimoram o desenho. Cada projeto possui
características próprias, mas um repertório consagrado de técnicas orienta essas
transformações; a seguir, destacam-se as mais pertinentes ao contexto deste trabalho.

\subsection{Composição de Métodos}

Um primeiro grupo de técnicas organiza o corpo das funções. \emph{Extrair Método}
(\emph{Extract Method}) transforma um fragmento de código coeso em um método próprio, com
nome descritivo, elevando a legibilidade e favorecendo o reúso; \emph{Embutir Método}
(\emph{Inline Method}) faz o inverso, eliminando indireções desnecessárias quando o corpo
de um método é tão evidente quanto o seu nome. De forma análoga, \emph{Extrair Variável}
(\emph{Extract Variable}) atribui a uma variável bem nomeada o resultado de uma expressão
complexa, tornando explícita a sua intenção.

\subsection{Movimentação de Recursos entre Objetos}

Quando as responsabilidades estão mal distribuídas, técnicas de movimentação reequilibram o
desenho. \emph{Mover Método/Campo} (\emph{Move Method/Field}) desloca um membro para a
classe que dele mais se beneficia; \emph{Extrair Classe} (\emph{Extract Class}) cria uma
nova classe quando uma única passa a acumular responsabilidades que seriam de duas. Esse
princípio de agrupar responsabilidades afins é diretamente análogo à forma como o ROS
organiza a computação em \emph{nós} especializados, cada qual responsável por uma parte bem
delimitada do sistema.

\subsection{Simplificação de Expressões Condicionais}

A complexidade condicional é uma fonte frequente de defeitos. \emph{Decompor Condicional}
(\emph{Decompose Conditional}) extrai a condição e as ações de ramificações complexas para
métodos nomeados; \emph{Consolidar Expressão Condicional} unifica verificações sucessivas
que conduzem ao mesmo resultado; \emph{Substituir Condicional por Polimorfismo} troca
seleções por tipo (por exemplo, um \texttt{switch}) por herança e polimorfismo; e as
\emph{Cláusulas de Guarda} (\emph{Guard Clauses}) tratam casos excepcionais com retornos
antecipados, mantendo limpo o ``caminho feliz'' da função. Em Python --- linguagem
predominante nos nós do EMA ---, tais técnicas reduzem o aninhamento e o excesso de
condicionais, alinhando-se à ênfase da linguagem em legibilidade.

\subsection{Organização de Dados}

Por fim, técnicas de organização de dados aumentam a clareza e a segurança do código.
\emph{Substituir Número Mágico por Constante} elimina valores literais soltos em favor de
constantes descritivas, e \emph{Encapsular Campo} restringe o acesso direto a atributos por
meio de métodos, controlando as invariantes do objeto. Em sistemas de reabilitação, nos
quais parâmetros como limiares de corrente e faixas angulares têm significado clínico
preciso, nomear e encapsular esses valores é também uma medida de segurança.

\subsection{Quando (e Quando Não) Refatorar}

A refatoração é mais eficaz quando integrada à rotina de desenvolvimento, e não tratada
como um evento isolado. A \emph{Regra do Escoteiro} sintetiza essa postura: deixar o código
mais limpo do que se o encontrou. Oportunidades naturais surgem ao adicionar
funcionalidades, ao corrigir defeitos e durante revisões de código. Um pressuposto
importante, contudo, é a existência de uma rede de segurança que garanta a preservação do
comportamento: sem testes automatizados, refatorar aproxima-se de ``mudar o código e
torcer''. O ciclo de Desenvolvimento Orientado a Testes (\emph{Test-Driven Development},
TDD) --- escrever um teste que falha, fazê-lo passar com o código mais simples possível e,
então, refatorar \cite{beck2003} --- oferece essa rede. Há, ainda, situações em que
refatorar é desaconselhável: nas proximidades de um prazo, pelo risco de introduzir
defeitos; e quando o código está tão comprometido, ou desprovido de testes, que a
reescrita se torna preferível --- afinal, a refatoração pressupõe que o sistema atual
funcione.

\section{Documentação Técnica}
\label{sec:documentacao}

Para os sistemas em estudo, a documentação é a maior carência isolada. Historicamente,
adotou-se no laboratório a premissa implícita de que ``o código é a documentação'' ---
premissa que se revela insuficiente quando um novo integrante precisa instalar, executar e
compreender o sistema a partir do zero. Esta seção apresenta os arcabouços que orientam uma
documentação intencional e sustentável.

\subsection{Arquitetura de Conteúdo: o Arcabouço Diátaxis}

O arcabouço \emph{Diátaxis} organiza a documentação técnica em quatro categorias definidas
pela necessidade do usuário \cite{diataxis}: \emph{tutoriais}, orientados ao aprendizado,
que conduzem o iniciante por um percurso passo a passo; \emph{guias de como fazer}
(\emph{how-to guides}), orientados a resolver tarefas específicas; \emph{referência},
orientada à informação precisa (parâmetros, APIs, formatos de dados); e
\emph{explicação}, orientada à compreensão dos conceitos e das decisões de projeto. A
distinção é útil porque cada categoria falha quando mistura propósitos --- um tutorial
sobrecarregado de detalhes de referência, por exemplo, deixa de cumprir o seu papel
didático. Para um laboratório com alta rotatividade, os tutoriais e os guias de instalação
são precisamente os artefatos cuja ausência mais onera a integração de novos membros.

\subsection{Documentação de Arquitetura: Modelo C4 e ADR}

A comunicação da arquitetura beneficia-se de notações padronizadas. O \emph{Modelo C4}
propõe descrever a arquitetura em quatro níveis de abstração progressivos --- Contexto,
Contêineres, Componentes e Código ---, permitindo ``dar zoom'' do panorama geral ao detalhe
conforme a audiência \cite{c4model}. Complementarmente, os \emph{Registros de Decisão de
Arquitetura} (\emph{Architecture Decision Records}, ADR) documentam, de forma incremental e
datada, as decisões arquiteturais relevantes, o contexto em que foram tomadas e as suas
consequências \cite{adr}. Os ADR são especialmente valiosos em projetos de longa duração,
pois preservam o \emph{porquê} das escolhas --- conhecimento que, de outro modo, se perde a
cada transição de equipe.

\subsection{Documentação como Código e Documentação Viva}

Duas abordagens operacionais aumentam a probabilidade de a documentação permanecer
atualizada. A \emph{documentação como código} (\emph{docs-as-code}) trata os artefatos de
documentação com o mesmo rigor do código-fonte: escritos em linguagens de marcação leve
(como Markdown), versionados em Git, submetidos a revisão por pares e publicados por
pipelines automatizados. A \emph{documentação viva} (\emph{living documentation}) gera parte
da documentação diretamente a partir do código ou de testes executáveis, reduzindo o risco
de defasagem. Ambas convergem para um mesmo objetivo: fazer da documentação um subproduto
natural do fluxo de trabalho, e não uma tarefa avulsa e facilmente adiada. Incorporar a
atualização da documentação à \emph{Definição de Pronto} (\emph{Definition of Done}) de
cada tarefa é uma prática de governança que reforça essa cultura.

\section{Versionamento e Organização de Repositórios}
\label{sec:versionamento}

O sistema de controle de versão é, em muitos projetos acadêmicos, a principal --- por vezes
a única --- memória institucional do software. Ferramentas como o Git, e plataformas como o
GitHub, oferecem recursos que vão muito além do simples armazenamento de arquivos:
organização por equipes, rastreamento de tarefas e defeitos (\emph{issues}), fluxos de
revisão por \emph{pull request}, marcação de versões (\emph{releases}) e integração
contínua. O uso disciplinado desses recursos --- convenções de nomenclatura de
repositórios, presença de arquivos \texttt{README} e de licença, histórico de \emph{commits}
inteligível e separação clara entre código ativo e arquivado --- é o que permite que o
conhecimento produzido seja localizável e reaproveitável.

A ausência dessas convenções produz o quadro oposto: repositórios duplicados para o mesmo
projeto, nomes que não comunicam propósito, misturas entre trabalhos originais e cópias, e
código sem qualquer ponto de entrada documentado. Como se detalha no
Capítulo~\ref{cap:diagnostico}, esse é exatamente o cenário observado na organização do
laboratório, e constitui uma das lacunas de Engenharia de Software mais custosas ---
justamente por ser das mais simples de corrigir com boas práticas.

\section{Arquitetura Baseada em Componentes e o ROS}
\label{sec:ros}

A Engenharia de Software Baseada em Componentes (\emph{Component-Based Software
Engineering}, CBSE) propõe construir sistemas pela composição de partes independentes e
substituíveis, que interagem por meio de interfaces bem definidas. O \emph{Robot Operating
System} (ROS) é a materialização desse paradigma no domínio da robótica: não um sistema
operacional no sentido estrito, mas um \emph{middleware} que organiza a computação em
\emph{nós} --- processos especializados e independentes --- que se comunicam por um modelo
de \emph{publicação/assinatura} (\emph{publish/subscribe}) sobre \emph{tópicos}
nomeados \cite{ros2docs}. Essa arquitetura favorece diretamente o baixo acoplamento: um nó
pode ser desenvolvido, testado, substituído ou reescrito sem afetar os demais, desde que
preserve o contrato das mensagens que troca.

O ROS~2, adotado como alvo pelo laboratório, assenta a sua comunicação sobre o padrão
\emph{Data Distribution Service} (DDS). Dessa escolha decorrem duas propriedades relevantes
para sistemas de reabilitação. A primeira é a \emph{descoberta automática} de nós: sem a
necessidade de um nó mestre central, os participantes localizam-se na rede e passam a
comunicar-se de forma descentralizada. A segunda são as políticas de \emph{Qualidade de
Serviço} (\emph{Quality of Service}, QoS), que permitem ajustar a comunicação à criticidade
de cada fluxo de dados. Assim, a telemetria de um sensor inercial, contínua e tolerante a
perdas, pode ser configurada como \emph{best effort} (privilegiando a velocidade), ao passo
que um comando enviado ao estimulador elétrico, que não pode ser perdido, é configurado
como \emph{reliable} (garantindo a entrega). Sobre esses mesmos mecanismos é possível
construir salvaguardas de segurança --- por exemplo, associar um tempo de validade
(\emph{lifespan}) e sinais de vivacidade (\emph{liveliness}) aos comandos críticos, de modo
que o driver do estimulador entre em modo seguro caso deixe de receber comandos válidos do
nó de controle. Esses conceitos são retomados na análise do estudo de caso do ciclismo com
eletroestimulação (Capítulo~\ref{cap:caso-trike}).

\section{Fundamentos de Domínio}
\label{sec:dominio}

Delimitados os fundamentos de Engenharia de Software, esta seção apresenta o domínio de
aplicação na extensão necessária para compreender os dois estudos de caso: a reabilitação
de marcha e a neuroplasticidade, a realidade virtual e a sua aplicação terapêutica, a
eletroestimulação funcional aplicada ao ciclismo e a captura de movimento por sensores
inerciais.

\subsection{Reabilitação de Marcha e Neuroplasticidade}
\label{sec:neuroplasticidade}

A marcha humana é uma tarefa motora complexa, resultante da coordenação fina entre o
sistema nervoso central, o sistema musculoesquelético e a integração contínua de
informações sensoriais. Lesões neurológicas como o AVC e a lesão medular interrompem, em
graus variados, o controle dessa tarefa. A reabilitação de marcha busca restaurar essa
capacidade explorando um princípio fundamental do sistema nervoso: a sua plasticidade.
Entende-se por \emph{neuroplasticidade} a capacidade do sistema nervoso de reorganizar a
sua estrutura e a sua função em resposta à experiência. \citeonline{kleim2008}
sistematizam esse conhecimento em princípios da plasticidade dependente de experiência,
entre os quais a especificidade do treino, a importância da repetição e da intensidade e a
relevância da saliência do estímulo. Esses princípios explicam \emph{por que} um treino
motor intensivo, repetitivo e engajador promove recuperação funcional
\cite{takeuchi2013}. É justamente sobre esses fatores --- repetição, intensidade e
saliência --- que operam as tecnologias assistivas discutidas a seguir
\cite{rajashekar2024,macchitella2023}.

\subsection{Realidade Virtual, Aumentada e Mista}
\label{sec:rv-ra-rm}

A realidade virtual (RV) pode ser definida como o uso de computação gráfica interativa para
gerar ambientes tridimensionais nos quais o usuário percebe estar imerso e com os quais
interage em tempo real. Dois conceitos caracterizam essa experiência: segundo
\citeonline{slater1997}, a \emph{imersão} é uma propriedade objetiva do sistema --- o grau
em que os dispositivos substituem os estímulos reais por estímulos sintéticos ---, ao passo
que a \emph{presença} é o estado subjetivo resultante, a sensação de \emph{estar} no
ambiente virtual. Para situar a RV em relação a tecnologias correlatas, é clássica a
taxonomia de \citeonline{milgram1994}, que organiza as combinações entre elementos reais e
virtuais ao longo de um \emph{continuum realidade--virtualidade}, ilustrado na
Figura~\ref{fig:continuum}.

\begin{figure}[H]
	\centering
	\caption{Continuum realidade--virtualidade de Milgram e Kishino.}
	\label{fig:continuum}
	\vspace{6pt}
	\begin{tikzpicture}[font=\small]
		\draw[{Latex[length=3mm]}-{Latex[length=3mm]}, thick] (0,0) -- (13,0);
		\foreach \x in {0,4.33,8.67,13}{\draw[thick] (\x,0.12) -- (\x,-0.12);}
		\node[align=center, anchor=north] at (0,-0.25) {Ambiente\\real};
		\node[align=center, anchor=north] at (4.33,-0.25) {Realidade\\aumentada\\(RA)};
		\node[align=center, anchor=north] at (8.67,-0.25) {Virtualidade\\aumentada\\(VA)};
		\node[align=center, anchor=north] at (13,-0.25) {Ambiente\\virtual};
		\draw[decorate, decoration={brace, amplitude=6pt}] (0.15,0.55) -- (12.85,0.55)
			node[midway, above=8pt]{Continuum realidade--virtualidade};
		\draw[decorate, decoration={brace, mirror, amplitude=6pt}] (4.33,-1.75) -- (8.67,-1.75)
			node[midway, below=8pt]{Realidade mista (RM)};
	\end{tikzpicture}

	\vspace{4pt}
	{\footnotesize Fonte: elaborado pelos autores com base em \citeonline{milgram1994}.}
\end{figure}

Embora concebida antes da difusão dos dispositivos atuais, essa taxonomia permanece a
referência conceitual dominante, ainda que trabalhos recentes proponham revisitá-la à luz
da evolução tecnológica \cite{skarbez2021}. Para os fins deste trabalho, o continuum de
\citeonline{milgram1994} é suficiente para posicionar o sistema em estudo: um ambiente
virtual imersivo, no extremo virtual do continuum, no qual o movimento real do paciente é
capturado e reproduzido.

\subsection{Realidade Virtual Aplicada à Reabilitação}
\label{sec:rv-reabilitacao}

A aplicação da RV à reabilitação motora fundamenta-se na possibilidade de oferecer, em
ambiente seguro e controlado, prática intensiva e repetitiva associada a realimentação
sensorial enriquecida --- precisamente os fatores associados à neuroplasticidade. A
magnitude desse benefício deve ser avaliada com rigor: a revisão sistemática Cochrane de
\citeonline{laver2017} conclui que a RV não é superior à terapia convencional quando
empregada isoladamente, mas mostra-se benéfica como \emph{adjuvante}, com efeitos mais
consistentes sob maior dose de treino. Um aspecto determinante é a forma como o sistema
apresenta ao usuário o seu próprio movimento --- a \emph{visualização do movimento} ---,
cujas principais formas, identificadas por \citeonline{ferreiradossantos2016}, a
Figura~\ref{fig:visualizacao} sintetiza.

\begin{figure}[H]
	\centering
	\caption{Formas de visualização do movimento em realidade virtual.}
	\label{fig:visualizacao}
	\vspace{6pt}
	\begin{tikzpicture}[
		font=\small,
		box/.style={draw, rounded corners, align=center, inner sep=5pt, minimum height=1cm},
		raiz/.style={draw, rounded corners, align=center, inner sep=6pt, fill=black!8, minimum height=1cm},
		]
		\node[raiz] (r) {Visualização do movimento\\em realidade virtual};
		\node[box, below=1.2cm of r, xshift=-5cm] (a) {Avatar\\(1\textsuperscript{a}/3\textsuperscript{a} pessoa)};
		\node[box, below=1.2cm of r, xshift=-1.7cm] (b) {Espelho\\virtual};
		\node[box, below=1.2cm of r, xshift=1.7cm] (c) {Realidade\\aumentada};
		\node[box, below=1.2cm of r, xshift=5cm] (d) {Representação\\abstrata/simbólica};
		\foreach \n in {a,b,c,d}{\draw[-{Latex[length=2mm]}] (r.south) -- (\n.north);}
	\end{tikzpicture}

	\vspace{4pt}
	{\footnotesize Fonte: elaborado pelos autores com base em \citeonline{ferreiradossantos2016}.}
\end{figure}

O sistema de RV do Projeto EMA, analisado no Capítulo~\ref{cap:caso-vr}, adota a estratégia
baseada em \emph{avatar}, reproduzindo em tempo real o movimento dos membros do paciente
capturado por sensores inerciais. A evidência sobre marcha e equilíbrio reforça o potencial
da abordagem quando associada a dispositivos de treino locomotor \cite{mirelman2016}, ainda
que os efeitos sobre parâmetros de marcha permaneçam inconsistentes \cite{lee2024}. Em
cenários de maior comprometimento, protocolos que associaram RV imersiva, realimentação
tátil e exoesqueletos controlados por interface cérebro-máquina chegaram a induzir
recuperação neurológica parcial em pacientes com lesão medular crônica \cite{donati2016}.

\subsection{Eletroestimulação Funcional e Ciclismo Assistido}
\label{sec:fes}

A eletroestimulação funcional (\emph{Functional Electrical Stimulation}, FES) consiste na
aplicação de pulsos elétricos a músculos paralisados ou enfraquecidos, provocando
contrações artificiais capazes de produzir movimento funcional. É a única técnica capaz de
restaurar movimento em músculos paralisados, permitindo, por exemplo, que indivíduos com
paraplegia executem o movimento de pedalada em uma bicicleta reclinada \cite{baptista2022}.
No ciclismo assistido por FES, sensores medem a posição angular dos pedais e, a partir
dela, determinam quais grupos musculares estimular e em que instante, de modo a produzir
uma pedalada coordenada. O grupo do LARA acumula experiência consolidada nesse tema,
inclusive com participação na competição internacional CYBATHLON, tendo desenvolvido desde
o software embarcado de controle do triciclo \cite{brindeiro2017} até aprimoramentos de
projeto centrados no usuário e técnicas de estimulação espacialmente distribuídas
\cite{baptista2022}. Esse sistema constitui o segundo estudo de caso deste trabalho.

\subsection{Captura de Movimento e Sensoriamento Inercial}
\label{sec:imu}

Ambos os estudos de caso dependem da captura do movimento do paciente por
\emph{unidades de medição inercial} (\emph{Inertial Measurement Units}, IMU) --- sensores
que combinam acelerômetros, giroscópios e, frequentemente, magnetômetros para estimar a
orientação de um segmento corporal ou de um componente mecânico no espaço. No sistema de
RV, as IMUs fixadas aos membros do paciente --- superiores e inferiores --- fornecem as
orientações que animam, em tempo real, os membros do avatar; no sistema de ciclismo, as IMUs
estimam a posição angular dos pedais que orienta a máquina de estados da estimulação. Em
ambos os casos, o
sensoriamento inercial é a ponte entre o movimento físico real e a lógica de software ---
razão pela qual a integridade e a documentação dos seus drivers são críticas para a
reprodução dos sistemas.
